%%====================================================================
\section{Conclusions}
\label{sec:conclusions}

We have presented \textcolor{red}{\textbf{tool name}}, a validated desktop application for gene-disease association extraction from full-text biomedical literature. The hybrid NER-LLM architecture---combining PubTator3's precision with Gemini Flash's recall---achieves zero hallucinated gene symbols across benchmark runs while maintaining competitive F1 scores (0.668 for cancer genomics, 0.611 for GWAS). The multi-layer validation framework (HGNC verification, grounding check, variant pattern matching, citation cross-referencing) provides defence-in-depth against LLM hallucinations. A controlled 36-run experiment demonstrated that multimodal figure analysis improves precision by $\Delta$F1\,=\,+0.833 on text-rich genomics papers.

The software is freely available under the MIT licence as a cross-platform desktop application, requiring only a free Google Gemini API key to operate. At \$0.10 per paper, it enables individual researchers to perform literature-scale extraction that would otherwise require teams of curators.

\paragraph{Future work.} Planned enhancements include: (1)~expanded variant extraction coverage for complex structural variants and non-HGVS notation---current coverage handles 12 common notation classes, but structural rearrangements remain an open challenge; (2)~integration with clinical databases (ClinVar, OMIM) for automated evidence cross-referencing, allowing the tool to flag whether an extracted gene--disease association has existing clinical evidence and directly supporting VUS interpretation workflows; (3)~inter-rater reliability validation of the gold standard, as the current benchmark uses single-annotator labels and Cohen's~$\kappa$ with an independent genetics expert is required before precision/recall metrics can be formally reported; (4)~specialised handling of immunoglobulin and HLA gene nomenclature, whose IMGT and IPD conventions conflict with HGNC identifiers and cause false negatives in immunogenomics papers; (5)~a pharmacogenomics-specific extraction mode targeting dosing recommendations and clinical action categories rather than molecular findings; and (6)~multi-paper relationship inference to aggregate gene--disease evidence across a corpus and identify consensus associations, as currently each paper is processed independently.
